\scanpage{37}
% archive sheet number: 36
\begingroup\makeatletter\tagsleft@true\makeatother
\uline{Cor} \deleted{D}ans le cas où $\operatorname{Card}\{\ \}=4$, \inserted{et $\Sigma_1^-\ne X,\Sigma_2^-\ne X$} les seules relations d'inclusions entre les ensembles $\Sigma_1^+,\Sigma_1^-,\Sigma_2^+,\Sigma_2^-$ sont les relations
\margin{Sous forme : Pour que les relations \unsure{strictes} $(*)$ soient les seules relations d'inclusion entre les ensembles $\Sigma_1^+,\Sigma_1^-,\Sigma_2^+,\Sigma_{\unsure{1|2}}^-$, il f et s. que ceux-ci soient \inserted{\unsure{mutlt.|mult.}} distincts et que l'on ait $\Sigma_1^-\ne X,\Sigma_2^-\ne X$.}
\[
\left\{\begin{aligned}
\Sigma_1^{\unsure{\top|+}}&\subset\Sigma_2^-\\
\Sigma_2^+&\subset\Sigma_1^-
\end{aligned}\right.\qquad .
\tag*{($*$)}
\]
En effet, on n'a pas les inclusions inverses, par hypothèse. D'autre part, on n'\uline{a pas}
\[
\Sigma_1^+\subset\Sigma_1^-\quad\text{ou}\quad
\Sigma_1^-\unsure{\subset|\in}\Sigma_1^+,\quad\text{car on aurait}
\]
$\deleted{\Sigma_1^+,\unsure{\Sigma_1|\Sigma_1^+|\Sigma_1^-}}\ \Sigma_1^-=X$, de $=$ on n'a pas $\Sigma_2^+\subset\unsure{\deleted{\vert}}{\unsure{\Sigma|E}}_{\unsure{1|2}}^-$ ni $\Sigma_2^-\subset\Sigma_2^+$, et on n'a pas $\Sigma_1^+\subset\Sigma_2^+$ ni l'inclusion inverse, puisque $\Sigma_1^+$ et $\unsure{\Sigma|\Sigma_2}^{\top}$ sont disjoints et $\ne\varnothing$. Il reste à voir que l'on n'a pas une inclusion $\Sigma_1^-\subset\Sigma_2^-$ ou inclusion inverse, on lit sur le diagramme $(14)$ que \deleted{\unsure{cela}} \inserted{la première \unsure{p.ex.|p. ex.}} signifierait que $\deleted{\Sigma_{12}\simeq\Sigma_1^-}$ \deleted{\unsure{d'où|de}} $\Sigma_2^-=X$, ce que l'on a exclu.

% The two thin, slightly wavering segments below reproduce the left marginal rule.
\noindent\begin{minipage}[t]{.04\linewidth}
\vspace{1.2\baselineskip}
\figure{\begin{tikzpicture}[x=.3pt,y=.3pt]
\draw[line width=.35pt] (0,0) -- (1,-75) -- (4,-155) -- (6,-212) -- (4,-270) -- (9,-310);
\draw[line width=.35pt] (9,-321) -- (13,-360) -- (7,-415) -- (5,-475) -- (1,-533) -- (3,-567);
\end{tikzpicture}}
\end{minipage}\hfill\begin{minipage}[t]{.94\linewidth}
\vspace{0pt}
\uline{\unsure{Dêf|Déf}} On dit que $\unsure{\Pi|\pi}_1,\unsure{\Pi|R}_2$ sont \uline{\unsure{strictement|Strictement} non\unsure{:|.} panachées} \inserted{si elles sont non panachées, et} si l\unsure{\deleted{es}}es quatre ensembles $\Sigma_1^\top,\Sigma_1^-,\Sigma_2^+,\Sigma_2^-$ sont distincts, \unsure{et.|i.e.} si \unsure{((|c} de plus $\Sigma_1^-\ne X,\Sigma_2^-\ne X$ -- i.e si les seules relations d'inclusion entre les ensembles $\Sigma_1^\top,\Sigma_1^-,\Sigma_2^\top,\Sigma_2^-$ sont $\Sigma_1^+\subset\Sigma_2^-,\Sigma_2^+\subset\Sigma_1^-$ -- ce qui signifie aussi que l'on a :
\[
\Sigma_1^-\ne X,\Sigma_2^-\ne X
\quad(\text{on dit donc que }\Pi_{\unsure{1|2}},\Pi_2\text{ sont \uline{propres}}),
\ E_1\ne\Sigma_{12},E_2\ne\Sigma_{12}.
\tag*{(21)}
\]
\end{minipage}
\endgroup
\clearpage
